\SourcePage{460}{463}
\section{Irreducible representations and the classification of terms}

Applications of group theory in quantum mechanics rest on the invariance of the Schr\"odinger equation for a physical system (an atom or molecule) under that system's symmetry transformations\footnote{Group-theoretical methods were first introduced into quantum mechanics by Wigner (E.~P.~Wigner, 1926).}. It follows directly that applying group elements to a function which solves the Schr\"odinger equation for some energy eigenvalue must produce further solutions of the same equation with that same energy. In other words, a symmetry transformation mixes among themselves the wave functions of the system's stationary states belonging to one energy level; these functions therefore realize a representation of the group. The essential point is that this representation is irreducible. Indeed, functions that necessarily mix with one another under symmetry transformations must in any event belong to the same energy level. By contrast, if the basis of a reducible%
\SourcePage{461}{464}%
representation were separated into several sets of functions that do not mix with one another, equality of the corresponding energy eigenvalues would be an exceedingly unlikely coincidence\footnote{Unless special reasons exist for it. In this connection, recall the ``accidental'' degeneracy that arises when the system's Hamiltonian has a symmetry higher than the purely geometrical symmetry discussed in this chapter (cf. the end of \S~36).}.

Thus each energy level of a system is associated with an irreducible representation of its symmetry group. The representation's dimension gives the degree of degeneracy of the level, namely the number of distinct states having that energy. Specifying the irreducible representation fixes all symmetry properties of the state---its response to the various symmetry transformations.

Irreducible representations of dimension greater than one occur only for groups containing elements that do not commute (Abelian groups possess only one-dimensional irreducible representations). It is worth recalling here that the connection between degeneracy and the existence of mutually noncommuting operators which nevertheless commute with the Hamiltonian was established earlier by arguments not involving group theory (see \S~10).

An important qualification must be attached to all these statements. As was pointed out previously (see \S~18), in the absence of a magnetic field, symmetry under reversal of the sign of time entails in quantum mechanics that complex-conjugate wave functions belong to the same energy eigenvalue. Consequently, when a set of functions and its complex-conjugate set realize two distinct (inequivalent) irreducible representations of the group, these conjugate representations must be taken together as one ``physically irreducible'' representation of twice the dimension; this convention will be understood everywhere below. Examples occurred in the preceding section. The group $C_3$, for instance, has only one-dimensional representations, yet two of them are complex conjugates and physically correspond to doubly degenerate energy levels. (In a magnetic field there is no symmetry under reversal of the sign of time, and complex-conjugate representations therefore correspond to different energy levels\footnote{Strictly, real characters (that is, equivalence of complex-conjugate representations) do not by themselves suffice%
\SourcePage{462}{465}%
to ensure that real basis functions can be chosen for a group representation. They do suffice for irreducible representations of point groups, although not for ``double'' point groups---see \S~99.}).

Suppose a physical system is acted on by a perturbation (the system is placed in an external field). To what extent can this perturbation split degenerate levels? The external field has a symmetry of its own\footnote{One example is provided by the energy levels of ionic $d$ and $f$ shells in a crystal lattice, weakly interacting with the surrounding atoms. Here the perturbation (external field) is the field exerted on the ion by the other atoms.}. If that symmetry is the same as, or higher than, the symmetry of the unperturbed system\footnote{If the symmetry group $H$ is a subgroup of $G$, the symmetry associated with $H$ is said to be lower than the higher symmetry of $G$. Plainly, for a sum of two expressions having symmetries $G$ and $H$, respectively, the symmetry is the lower symmetry $H$.}, then the perturbed Hamiltonian $\widehat H=\widehat H_0+\widehat V$ has the same symmetry as the unperturbed operator $\widehat H_0$. No splitting of degenerate levels can then occur. If, however, the perturbation has lower symmetry than the unperturbed system, the symmetry of the Hamiltonian $\widehat H$ will be that of the perturbation $\widehat V$. The wave functions which formed an irreducible representation of the symmetry group of $\widehat H_0$ will still form a representation of the symmetry group of the perturbed operator $\widehat H$, but this representation may be reducible; the degenerate level then splits. The following example shows how the machinery of group theory gives a definite answer concerning the splitting of a particular level.

Let the unperturbed system have symmetry $T_d$. Consider a triply degenerate level corresponding to the irreducible representation $F_2$ of this group. Its characters are
\[
\begin{array}{c|ccccc}
 & E&8C_3&3C_2&6\sigma_d&6S_4\\ \hline
F_2&3&0&-1&1&-1
\end{array}
\]
Now let a perturbation of symmetry $C_{3v}$ act on the system, with its threefold axis coinciding with one of the corresponding axes of $T_d$. The three wave functions of the degenerate level realize a representation of $C_{3v}$,%
\SourcePage{463}{466}%
a subgroup of $T_d$. The characters of this representation are simply those of the same elements in the original $T_d$ representation:
\[
\begin{array}{c|ccc}
 &E&2C_3&3\sigma_v\\ \hline
 &3&0&1
\end{array}
\]
This representation, however, is reducible. From the characters of the irreducible representations of $C_{3v}$, it is readily decomposed into irreducible parts by the general rule (94.16). The result is the sum of the $A_1$ and $E$ representations of $C_{3v}$. Hence the triply degenerate $F_2$ level separates into one nondegenerate $A_1$ level and one doubly degenerate $E$ level. If instead the same system is subjected to a perturbation of symmetry $C_{2v}$, also a subgroup of $T_d$, the wave functions of the same $F_2$ level give a representation whose characters are
\[
\begin{array}{c|cccc}
 &E&C_2&\sigma_v&\sigma'_v\\ \hline
 &3&-1&1&1
\end{array}
\]
Decomposition into irreducible parts shows that it contains the representations $A_1$, $B_1$, and $B_2$. In this case, therefore, the level splits completely into three nondegenerate levels.
